\documentclass{article}

% Workshop track (double-blind) — most appropriate for a reference
% architecture and validation harness without trained-model results.
\usepackage[dblblindworkshop]{neurips_2026}

\usepackage[utf8]{inputenc}
\usepackage[T1]{fontenc}
\usepackage{hyperref}
\usepackage{url}
\usepackage{booktabs}
\usepackage{amsfonts}
\usepackage{amsmath}
\usepackage{amssymb}
\usepackage{nicefrac}
\usepackage{microtype}
\usepackage{xcolor}
\usepackage{graphicx}
\usepackage{array}
\usepackage{multirow}

\title{A Reference Architecture and Validation Harness for Multimodal Emotion Recognition in Mental Health Diagnostics}

\author{%
  Anonymous Author(s) \\
  Affiliation(s) \\
  \texttt{email(s)} \\
}

\begin{document}

\maketitle

\begin{abstract}
Multimodal emotion recognition (MER) offers a promising avenue for early detection and
monitoring of mental health conditions by analysing speech prosody, facial expressions,
and textual content from clinical interviews. Despite advances in unimodal encoders and
multimodal fusion, existing MER systems have not been validated for clinical deployment:
they lack robustness to clinically-correlated modality dropout, ignore the session-level
temporal structure of therapeutic interviews, and provide no interpretability grounded in
diagnostic constructs such as the Patient Health Questionnaire (PHQ-8). We present a
reference architecture---the Cross-Modal Gated Exchange Fusion network with Session-Level
Hierarchical Transformer and Clinical Concept Bottleneck (CM-GEF + SLHT + CCB)---together
with a comprehensive validation harness. The architecture introduces per-dimension gated
feature substitution informed by learned modality quality estimates, therapeutic phase
embeddings that contextualise utterances within a 4-phase interview structure, and a
concept bottleneck that maps session representations through DSM-5-aligned clinical
constructs. The harness comprises 35+ unit tests covering shape correctness, gradient flow,
and bf16 numerical stability; seven synthetic benchmarks for missing-modality stress,
clinical dropout patterns, and sustained affect detection; and ten single-field ablations
that isolate each architectural decision. We release the typed \texttt{ModelConfig}
contract, PyTorch reference implementation, and validation suite as an open-source
starting point for follow-up work on clinical MER. Trained-model evaluation on real
clinical datasets is explicitly left as future work.
\end{abstract}


\section{Introduction}

Mental health disorders affect an estimated one in eight people worldwide, with depression
alone contributing to over 50 million years lived with disability globally~\cite{who2023depression}.
Early detection and continuous monitoring of emotional states can substantially improve
therapeutic outcomes, yet many cases remain undiagnosed until they reach critical
stages. Multimodal emotion recognition (MER) addresses this gap by integrating cues from
speech, facial expressions, and language, offering a richer assessment of a patient's
emotional state than any single modality alone~\cite{survey2026multimodal}.

Recent years have seen significant progress in both unimodal emotion encoders and
multimodal fusion architectures. Self-supervised speech models such as WavLM~\cite{chen2022wavlm}
and HuBERT~\cite{hsu2021hubert} provide robust paralinguistic representations; vision
transformers like DINOv2~\cite{oquab2023dinov2} deliver strong facial affect features;
and domain-adapted language models such as MentalBERT~\cite{ji2022mentalbert} capture
clinical vocabulary. On the fusion side, cross-modal Transformers~\cite{tsai2019mult},
hierarchical gated networks~\cite{aefnet2025}, and graph-based approaches~\cite{micga2025}
have pushed benchmark accuracy on lab datasets to above 99\%~\cite{unifiedmer2025}.

Despite these advances, three critical gaps prevent existing MER systems from being
deployed in clinical mental health settings. \textbf{First}, missing-modality robustness
has only been evaluated under random dropout conditions on clean lab benchmarks
(CMU-MOSI, IEMOCAP)~\cite{amte2025,micga2025}. In clinical interviews, modality dropout
is systematically correlated with emotional state---patients turn away during distress
(vision lost), cry (speech unintelligible), or fall silent (text absent)---yet no
existing work validates gated fusion under such clinically-correlated patterns.
\textbf{Second}, the temporal structure of a 30--60 minute therapeutic interview---with
distinct phases of rapport-building, exploration, intervention, and closure---is ignored;
current models process short clips or windowed utterances~\cite{dsin2025,dstc2025}.
\textbf{Third}, model predictions are not interpretable in clinically actionable terms;
generated explanations~\cite{ecmc2026} are not grounded in standardised diagnostic
constructs such as the PHQ-8~\cite{kroenke2009phq}.

This paper presents a \textbf{reference architecture and validation harness} designed
to address these gaps. We do not claim to outperform existing systems on clinical
data---no model training on real clinical datasets has been conducted. Instead, our
contribution is the architecture itself, the design rationale substantiated by
inductive-bias arguments, and a reproducible validation infrastructure that enables
the community to iterate on clinical MER designs.

In summary, our contributions are:
\begin{itemize}
  \item We present a reference architecture for clinical MER that combines
        \textbf{Cross-Modal Gated Exchange Fusion} (CM-GEF) with learned per-utterance
        modality quality estimates, a \textbf{Session-Level Hierarchical Transformer}
        (SLHT) with therapeutic phase embeddings, and a \textbf{Clinical Concept
        Bottleneck} (CCB) aligned to PHQ-8 depression criteria.
  \item We provide a typed \texttt{ModelConfig} dataclass that encodes all 60+
        hyperparameters with defaults and docstring rationale, enabling systematic
        exploration of the design space.
  \item We release a validation harness covering shape correctness (35+ pytest tests),
        gradient flow verification, bf16 numerical stability, seven synthetic benchmarks
        (missing modality stress, clinical dropout patterns, late fusion comparison,
        quality estimator sensitivity, sustained affect proxy, concept mapping fidelity,
        throughput profiling), and ten single-field ablations.
  \item We open-source the full PyTorch reference implementation under a permissive
        license to support reproducible follow-up work.
\end{itemize}

The remainder of this paper is organised as follows. Section~\ref{sec:related} reviews
related work across unimodal emotion recognition, multimodal fusion, and clinical MER.
Section~\ref{sec:method} describes the proposed architecture in detail.
Section~\ref{sec:validation} presents the validation harness and experimental setup.
Section~\ref{sec:results} reports results from the harness evaluation.
Section~\ref{sec:discussion} discusses implications, design insights, and limitations.
Section~\ref{sec:conclusion} concludes with future research directions.


\section{Related Work}
\label{sec:related}

\subsection{Unimodal Emotion Recognition}

\textbf{Speech/Audio.} Self-supervised speech models have become the de facto standard for
paralinguistic emotion recognition. WavLM~\cite{chen2022wavlm} extends the HuBERT
framework with denoising and speaker-aware pretraining, achieving state-of-the-art
performance on emotion recognition benchmarks. Whisper encoder features have recently
shown promise as depression markers~\cite{sadeghi2024}. Conformer architectures combine
convolution and self-attention for streaming acoustic emotion recognition.
Despite strong performance on clean corpora, speech emotion recognition degrades sharply
under noise, codec compression, and non-English prosody, and cross-corpus generalisation
to clinical populations remains poor~\cite{survey2026multimodal}.

\textbf{Computer Vision.} Facial expression recognition (FER) has been transformed by
vision transformers. DINOv2~\cite{oquab2023dinov2} produces self-supervised features
that transfer strongly to facial affect tasks. Action Unit (AU) detection provides
anatomically grounded alternatives to holistic emotion classification and is less prone
to spurious correlations. Clinical FER faces additional challenges: variable lighting,
head pose, occlusions (masks), and limited labelled clinical video datasets.
Synthetic augmentation and domain adaptation for clinical FER remain
under-explored~\cite{survey2025mer}.

\textbf{Language Models.} For textual emotion and sentiment analysis, RoBERTa and
DeBERTa-v3 provide strong baselines. Domain-adapted variants such as
MentalBERT~\cite{ji2022mentalbert} and ClinicalBERT are fine-tuned on mental health
text and clinical notes, capturing specialised vocabulary. Large language models offer
zero-shot emotion reasoning with interpretable explanations but at high compute cost
and with a tendency to hallucinate emotion rationales. Clinical utterances are often
short, and inter-annotator agreement on clinical emotion labels is low~\cite{survey2025mer}.

\subsection{Multimodal Fusion Architectures}

Multimodal fusion is the central design challenge for MER. Approaches can be categorised
along a spectrum of cross-modal interaction depth. \textbf{Early fusion} concatenates
features from each modality before prediction; it is simple but ignores cross-modal
dynamics and requires all modalities at inference time. \textbf{Late fusion} ensembles
independent unimodal predictions, offering robustness to missing modalities but losing
fine-grained cross-modal interaction. \textbf{Intermediate fusion}---the current
state-of-the-art paradigm---includes cross-modal Transformers (MulT, 2019)~\cite{tsai2019mult},
which apply pairwise cross-attention between modalities; hierarchical gated fusion
(AEFNet, 2025)~\cite{aefnet2025}, which stacks unimodal, bimodal, and trimodal gates;
adaptive exchange fusion (AMTE, 2025)~\cite{amte2025}, which supplements weak local
features with strong global features from another modality; and graph-based fusion
(Mi-CGA, 2025)~\cite{micga2025}, which uses cross-modal graph attention for robustness
to missing modalities.

Unified frameworks combining variational mixture-of-experts with cross-attentional
modality interaction have recently achieved over 99\% accuracy on MELD and
IEMOCAP~\cite{unifiedmer2025}. However, these results are on curated lab datasets with
clean, well-aligned modalities; clinical deployment introduces noise, missing data,
and systematic dropout patterns that these frameworks have not been tested against.

\subsection{Clinical Mental Health MER}

A smaller but growing body of work applies MER specifically to depression detection.
MDD-MARF~\cite{zhou2025mddmarf} combines tri-modal (audio+video+text) multi-level
attention with residual fusion, achieving the current state-of-the-art on DAIC-WOZ with
MAE 3.13 and RMSE 3.59 on PHQ-8 prediction. The DSTC
framework~\cite{dstc2025} uses Temporal ProbSparse Attention with position-guided
multimodal cross-fusion, reporting MAE 4.22 on E-DAIC. Pr3+Whisper+AudioQual~\cite{sadeghi2024}
leverages LLM-based text features with speech quality filtering for fully automated
depression detection, achieving MAE 3.86 on E-DAIC.

Three critical gaps persist across all existing clinical MER work. \textbf{Missing-modality
robustness} has not been evaluated under clinical (non-random) dropout patterns---prior
studies ablate entire modalities but do not model the correlated dropout observed in
real therapy sessions. \textbf{Session-level temporal modelling} is absent: DSIN~\cite{dsin2025}
and DSTC~\cite{dstc2025} process short clips or aggregate via pooling, ignoring the
therapeutic phase structure of clinical interviews. \textbf{Clinical interpretability}
is lacking: ECMC~\cite{ecmc2026} generates free-text emotion captions but does not map
predictions to standardised diagnostic items. Our architecture is designed to address
all three gaps simultaneously.


\section{Proposed Architecture}
\label{sec:method}

Figure~\ref{fig:arch} provides an overview of the full architecture. The pipeline
proceeds through four stages: (1) unimodal encoding via pretrained backbones,
(2) cross-modal gated exchange fusion, (3) session-level hierarchical temporal modelling,
and (4) clinical concept bottleneck prediction.

\begin{figure}[h]
\centering
\begin{verbatim}
                        PHQ-8 Total (B,)
                             |
                        PHQ-8 Items (B, 8)
                             |
                   Clinical Concept Bottleneck
                   8 PHQ-8-aligned concepts (B, 8)
                             |
                   Session Representation (B, d_s)
                             |
           Session-Level Hierarchical Transformer
          Phase Emb. + Bidirectional MHSA × 4 + Pool
                             |
          Cross-Modal Gated Exchange Fusion (CM-GEF)
         Quality Est. + Exchange Attn + Gate + Aggregate
                             |
       +-------------+ +----------+ +----------------+
       | WavLM Base+ | | DINOv2   | | MentalBERT     |
       | (audio)     | | (video)  | | (text)         |
       | d=768       | | d=384    | | d=768          |
       +-------------+ +----------+ +----------------+
\end{verbatim}
\caption{Overview of the proposed MER architecture. Three pretrained unimodal encoders
feed into CM-GEF for cross-modal fusion. Fused utterance representations pass through
the SLHT for session-level temporal modelling with therapeutic phase awareness,
then through the CCB for clinically interpretable PHQ-8 prediction.}
\label{fig:arch}
\end{figure}

\subsection{Unimodal Encoders}

Each modality is processed by a pretrained encoder that remains frozen during Stage~1
fusion pretraining and receives LoRA adapters during Stage~2 clinical adaptation.
\textbf{Audio} is encoded by WavLM Base Plus~\cite{chen2022wavlm}, which processes 16~kHz
raw waveforms. Attention-weighted mean pooling produces a fixed-dimension utterance
representation of $d_{\text{audio}} = 768$. \textbf{Video} is encoded by DINOv2
Small~\cite{oquab2023dinov2}, which processes 30 uniformly sampled $224\times224$ facial
crops per utterance and produces $d_{\text{video}} = 384$ features via [CLS] token
pooling. \textbf{Text} is encoded by MentalBERT~\cite{ji2022mentalbert}, which tokenises
utterances up to 128 tokens and produces $d_{\text{text}} = 768$ features via [CLS] plus
attention-weighted sequence pooling. Each encoder feeds into a lightweight projection MLP
($d_m \to d_{\text{proj}} = 256$) that unifies all modalities to a common latent
dimension.

\subsection{Modality Quality Estimators}

A key design requirement is the ability to detect when a modality is unreliable. For each
utterance, a lightweight MLP head predicts the modality's quality from the pooled encoder
features~\cite{amte2025}. The quality estimator is a three-class classifier---clean,
degraded, or absent---trained on synthetic corruptions (noise injection, face occlusion,
text truncation). The reliability score $r_m \in [0, 1]$ for modality $m$ is the
softmax probability of the ``clean'' class:

\begin{equation}
  \mathbf{h} = \text{GELU}(\mathbf{W}_1 \text{LayerNorm}(\mathbf{f}_m)),
  \quad
  \mathbf{p} = \text{softmax}(\mathbf{W}_2 \mathbf{h}),
  \quad
  r_m = p_0.
  \label{eq:quality}
\end{equation}

During Stage~2 clinical fine-tuning, the quality estimators can be run in inference mode
to prevent feature collapse from distribution shift. The reliability scores play a dual
role: they gate the cross-modal substitution and weight the final aggregation.

\subsection{Cross-Modal Gated Exchange Fusion (CM-GEF)}
\label{sec:cmgef}

CM-GEF is the central fusion mechanism. For each utterance, each modality $m$ computes a
cross-modal context vector by attending to the other two modalities via multi-head exchange
attention. A per-dimension gate then interpolates between the modality's own projected
features $\mathbf{f}_m$ and the cross-modal context $\mathbf{ctx}_m$, modulated by the
estimated reliability $r_m$:

\begin{align}
  \mathbf{ctx}_m &= \text{MultiHeadAttention}(\mathbf{Q}=\mathbf{f}_m,
                    \mathbf{KV}=[\mathbf{f}_j; \mathbf{f}_k]),
  \label{eq:exchange}\\
  \mathbf{gate}_m &= \sigma\big(\mathbf{W}_g[\mathbf{f}_m; \mathbf{ctx}_m; r_m]\big),
  \label{eq:gate}\\
  \mathbf{f}_m^{\text{fused}} &= \mathbf{gate}_m \odot \mathbf{f}_m
                               + (1 - \mathbf{gate}_m) \odot \mathbf{ctx}_m.
  \label{eq:substitution}
\end{align}

The per-dimension gating---rather than a scalar modality-level gate---is an intentional
design choice: a modality may be unreliable in some feature subspaces (e.g., audio timbre
distorted) while remaining reliable in others (e.g., audio pitch preserved).

The three gated representations are aggregated via softmax weights proportional to the
reliability scores, providing a second line of defence against unreliable modalities:

\begin{equation}
  \mathbf{f}^{\text{utt}} = \sum_{m \in \{a,v,t\}}
    \frac{\exp(r_m)}{\sum_{j}\exp(r_j)} \; \mathbf{f}_m^{\text{fused}}.
  \label{eq:aggregation}
\end{equation}

The exchange attention and gated substitution may be stacked for $L_{\text{ex}}$ layers
(default $L_{\text{ex}} = 2$) with residual connections and feed-forward blocks,
enabling iterative cross-modal refinement.

\paragraph{Design rationale.}
Per-dimension gating provides finer-grained control than scalar modality-level gates,
at the cost of increased capacity (which may overfit on small clinical datasets---we
mitigate this with gate entropy regularisation and dropout). Multi-head exchange
attention allows different heads to specialise on different cross-modal patterns
(e.g., attending to lip movements in video versus sentiment in text), a capability
well-established in multimodal Transformers~\cite{tsai2019mult}. The dual-gating
mechanism (gated substitution via Eq.~\ref{eq:substitution} plus reliability-weighted
aggregation via Eq.~\ref{eq:aggregation}) provides robustness redundancy; ablation
studies isolate the contribution of each.

\subsection{Session-Level Hierarchical Transformer (SLHT)}

Clinical interviews unfold over 30--60 minutes with a well-defined structure: rapport-building
(first~$\sim$5~min), exploration (open-ended probes), intervention (targeted clinical
questions), and closure (wrap-up). Depression markers manifest differently across these
phases---flat affect during rapport-building may indicate sustained depression, while
emotional reactivity during intervention probes is a positive prognostic sign.

SLHT captures this structure through four mechanisms:

\textbf{(1) Phase embedding.} Each utterance receives a learned embedding of its
therapeutic phase (one of four), which is concatenated with the fused utterance feature
before projection into the session model dimension.

\textbf{(2) Position encoding.} Learned position encodings index utterances by their
temporal order within the session (up to $T_{\max} = 512$).

\textbf{(3) Bidirectional Transformer.} $L_{\text{sess}} = 4$ layers of multi-head
self-attention and feed-forward blocks process the full utterance sequence. Bidirectional
(rather than causal) attention is used because the entire session is available for
post-hoc analysis, and early emotional cues gain context from later utterances.

\textbf{(4) Phase-level pooling.} For each of the four phases, utterances belonging to
that phase are pooled via cross-attention between a phase-level query and the full
utterance sequence, producing per-phase representations. A learned importance weight
over phases produces the final session representation:

\begin{equation}
  \mathbf{s} = \sum_{p=1}^{4} \alpha_p \, \mathbf{\psi}_p + \text{mean}(\mathbf{h}),
  \qquad
  \alpha_p = \text{softmax}(\mathbf{w}_{\text{imp}}^\top \mathbf{\psi}_p).
  \label{eq:session}
\end{equation}

The residual global context $\text{mean}(\mathbf{h})$ ensures that even if some phases
contain few utterances, the session representation retains a complete view.

\subsection{Clinical Concept Bottleneck (CCB)}

Rather than predicting PHQ-8 scores directly from session features, CCB first maps the
session representation through an interpretable intermediate layer of 8 clinical concepts
aligned to DSM-5 depression criteria (anhedonia, depressed mood, sleep disturbance,
fatigue, appetite change, guilt/worthlessness, concentration problems, psychomotor change).
Each concept score $c_k \in [0, 1]$ is produced by a sigmoid-activated MLP:

\begin{equation}
  \mathbf{c} = \sigma\big(\mathbf{W}_{c2}\,
               \text{GELU}(\mathbf{W}_{c1}\,
               \text{LayerNorm}(\mathbf{s}))\big).
  \label{eq:concept}
\end{equation}

PHQ-8 item scores $\mathbf{i} \in [0, 3]^8$ are computed from the concept scores via a
learned weight matrix with a diagonal-dominant prior (each concept $k$ primarily maps to
item $k$). The total score is the sum of items, and binary depression classification uses
a sigmoid over concept scores. This structure makes the model's reasoning transparent:
clinicians can inspect which concepts were activated and verify that the concept-to-item
mapping is clinically sensible, analogous to concept bottleneck models in other
domains~\cite{koh2020concept}.

\subsection{Two-Stage Training Pipeline}

The training pipeline follows a two-stage strategy designed for scarce clinical labelled
data. In \textbf{Stage~1} (fusion pretraining), the unimodal encoders are frozen while
the projection MLPs, quality estimators, CM-GEF, SLHT, and CCB are trained on large-scale
in-the-wild emotion data (CMU-MOSEI, MELD, IEMOCAP; approximately 47k multimodal
utterances) with a multi-task loss of emotion classification, modality dropout
reconstruction, and quality estimation. In \textbf{Stage~2} (clinical adaptation), LoRA
adapters ($r=8, \alpha=16$) are applied to the encoder backbones while the fusion and
prediction heads are fine-tuned on clinical data (DAIC-WOZ/E-DAIC; approximately 189--275
sessions) with a loss combining PHQ-8 MAE, binary cross-entropy for depression
classification, and optional concept supervision.


\section{Validation Infrastructure}
\label{sec:validation}

The validation harness is designed to verify architectural correctness and characterise
structural properties of the design. It operates entirely on synthetic data---no real
clinical datasets are required, and no model training has been conducted. All benchmarks
use randomly initialised weights to test the \emph{infrastructure}, not model performance.
Results are reported on a single CPU node (AMD EPYC 7B12, 8 cores) with bf16 tests
requiring an NVIDIA A100 80~GB GPU.

\subsection{Layer~1: Unit Tests}

The pytest suite contains 35+ tests organised into nine classes:

\begin{itemize}
  \item \textbf{Shape correctness} (9 tests): verifies that every component produces the
        expected output dimensions for single-utterance, full-session, and streaming
        modes, across variable utterance counts ($T \in \{4, 8, 16\}$).
  \item \textbf{Gradient flow} (7 tests): verifies that loss gradients reach all actively
        used trainable parameters, including CM-GEF, SLHT, CCB, quality estimators, and
        LoRA adapters (Stage~2). Checks for NaN and Inf gradients.
  \item \textbf{Numerical stability} (7 tests): verifies forward pass stability in
        bfloat16 and float16, under extreme input magnitudes ($\times 1000$), silent
        audio, and complete modality absence.
  \item \textbf{Domain correctness} (24 tests): covers missing modality robustness
        (single-modality input, modality dropout augmentation, weight normalisation),
        concept bottleneck range constraints (scores $\in [0,1]$, items $\in [0,3]$,
        total $\in [0,24]$), diagonal-dominant initialisation, quality estimator bounds,
        session transformer variable-length handling, and training pipeline stage
        configuration.
\end{itemize}

\subsection{Layer~2: Domain Benchmarks}

Seven synthetic benchmarks probe specific design properties:

\begin{enumerate}
  \item \textbf{Missing modality stress.} Evaluates under 0\%, 20\%, 40\%, and 60\%
        random dropout per modality, measuring MAE degradation from the full-modality
        baseline.
  \item \textbf{Clinical dropout patterns.} Three clinically-correlated patterns: vision
        drop during distress (40\% video missing), audio drop in noisy environment (30\%
        audio missing), and text drop when non-verbal (40\% text + 10\% audio/video).
  \item \textbf{Late fusion comparison.} Compares CM-GEF against a simple late fusion
        baseline (mean of available modality features) under 0\% and 40\% dropout.
  \item \textbf{Quality estimator sensitivity.} Measures whether reliability scores
        differ between clean and synthetically corrupted features.
  \item \textbf{Sustained affect proxy.} Compares session-level representation separation
        (SLHT) vs. utterance-level mean pooling for sessions with artificially induced
        ``flat affect'' segments (low-amplitude waveform).
  \item \textbf{Concept mapping fidelity.} Verifies diagonal dominance
        ($\text{diag}/\text{off-diag} > 2$) of the concept-to-item weight matrix.
  \item \textbf{Throughput profiling.} Measures inference latency, memory usage, and
        estimated FLOPs.
\end{enumerate}

\subsection{Layer~3: Ablations}

Ten single-field \texttt{ModelConfig} ablations isolate the contribution of each
architectural decision (Table~\ref{tab:ablations}). Each ablation compares a baseline
model against an otherwise identical model with one hyperparameter changed, measuring
MAE on clean and 20\%-dropout synthetic data.

\begin{table}[h]
\caption{Summary of the ten single-field ablations in the validation suite. Each
ablation tests a specific design hypothesis by modifying a single
\texttt{ModelConfig} field.}
\label{tab:ablations}
\centering
\begin{tabular}{lll}
\toprule
\# & Ablation & Config field change \\
\midrule
1 & Drop CM-GEF $\to$ late fusion & $n_{\text{exchange\_layers}}: 2 \to 0$ \\
2 & Scalar gates $\to$ per-dim gates & $\text{exchange\_hidden\_dim}: 128 \to 0$ \\
3 & Drop quality estimators & $n_{\text{audio\_quality\_classes}}: 3 \to 0$ \\
4 & Drop phase embeddings & $n_{\text{therapeutic\_phases}}: 4 \to 1$ \\
5 & Drop session Transformer $\to$ mean pool & $n_{\text{session\_layers}}: 4 \to 0$ \\
6 & LoRA vs. full fine-tune & $\text{lora\_r}: 8 \to 0$ \\
7 & Drop concept bottleneck & $n_{\text{clinical\_concepts}}: 8 \to 0$ \\
8 & Drop dropout simulation & $p_{\text{dropout}}: 0.15 \to 0.0$ \\
9 & WavLM Base+ vs. Small & $\text{audio\_encoder\_name}$ changed \\
10 & Bidirectional vs. causal session & $\text{streaming\_mode}: \text{False} \to \text{True}$ \\
\bottomrule
\end{tabular}
\end{table}


\section{Results}
\label{sec:results}

All results reported here are from the validation harness running on a randomly
initialised (untrained) model. They verify the correctness and structural properties of
the architecture---they are \emph{not} measures of clinical performance, which requires
trained-model evaluation on real clinical datasets.

\subsection{Architectural Correctness}

All 35+ unit tests pass across all test classes. Key structural verifications:

\textbf{Output shapes.} The model produces correct tensor shapes for all three operating
modes. Single-utterance encoding yields $(\text{batch}, d_{\text{proj}})$ with
$d_{\text{proj}} = 256$ at default config. Full-session forward yields
$(\text{batch},)$ for PHQ-8 total, $(\text{batch}, 8)$ for PHQ-8 items and concept
scores, $(\text{batch}, d_{\text{session}})$ for session representation, and
$(\text{batch}, 4, d_{\text{session}})$ for per-phase representations. Variable
session lengths $T \in \{4, 8, 16\}$ are handled correctly.

\textbf{Gradient flow.} All actively used trainable parameters receive non-NaN,
non-Inf gradients in both Stage~1 (fusion modules trainable, encoders frozen) and
Stage~2 (LoRA adapters trainable). Gradient flow reaches all four novel blocks:
CM-GEF, SLHT, CCB, and quality estimators.

\textbf{Numerical stability.} Forward passes are stable in bfloat16: PHQ-8 total and
item scores remain finite with no NaN or Inf values. Stability holds under extreme
input scaling ($\times 1000$), silent audio (all-zero waveforms), and complete modality
absence (all three modalities masked out). Gate values maintain diversity across the
batch (mean per-gate within $[0.05, 0.95]$), confirming the absence of gate saturation.

\subsection{Parameter Counts and Complexity}

The default configuration ($d_{\text{proj}} = 256$, $d_{\text{session}} = 256$,
$L_{\text{ex}} = 2$, $L_{\text{sess}} = 4$) yields approximately 23~M total parameters.
Of these, approximately 3~M are trainable in Stage~1 (projection MLPs, quality
estimators, fusion, session transformer, concept bottleneck) and approximately 10.5~M
are trainable in Stage~2 (adding LoRA adapters on encoder backbone projections).
Time complexity is $O(d_{\text{proj}}^2)$ per utterance for CM-GEF and
$O(T^2 \cdot d_{\text{session}})$ per session for the bidirectional SLHT, where
$T \leq 512$. The concept bottleneck contributes negligible overhead (approximately
0.5~M parameters).

\subsection{Synthetic Benchmark Results}

\paragraph{Concept mapping fidelity.} The concept-to-item weight matrix is initialised
with a diagonal-dominant prior. The ratio of diagonal to mean off-diagonal weight
magnitude exceeds 100 ($\gg 2$ pass criterion), confirming that each PHQ-8 item is
primarily influenced by its corresponding clinical concept at initialisation.

\paragraph{Quality estimator sensitivity.} With random initialisation, reliability
scores show limited sensitivity to synthetic corruption---differences between clean
and noisy inputs are below 0.001 on average. This is expected, as the quality estimators
require supervised training on synthetic corruptions before they become sensitive.
The training infrastructure (synthetic corruption labels, MLP classifier head,
multi-task loss) is in place and verified.

\paragraph{Throughput.} Profiling on CPU with the default config estimates approximately
4.4~GFLOP per sample (Kaplan~\etal{} scaling estimate for transformer-heavy models:
$6 \times \text{params} \times T$, with $T = 32$).
% TODO: not yet measured — requires GPU benchmark with trained model.

\subsection{Ablation Infrastructure}

All ten ablations execute successfully and produce quantitative comparisons
(clean MAE, dropout MAE, parameter count) between the baseline and ablated
configurations. The ablation runner verifies that each single-field change produces
a valid model without breaking the forward pass.
% TODO: not yet measured — effect sizes require trained-model evaluation on clinical data.


\section{Discussion}
\label{sec:discussion}

The validation harness confirms that the proposed architecture is structurally sound:
it produces correct tensor shapes across all operating modes, maintains gradient flow
through all trainable components, and remains numerically stable under extreme conditions
including bf16 precision, large input magnitudes, and complete modality absence. These
properties are necessary---though not sufficient---for clinical deployment.

\subsection{Design Insights from the Harness}

Several structural properties emerge from the harness characterisation. \textbf{The
triple-redundancy missing-modality strategy} (quality estimator $\to$ gated substitution
$\to$ reliability-weighted aggregation) is computationally lightweight: the quality
estimators add fewer than 0.15~M parameters, and the exchange attention is linear in
the number of source modalities. \textbf{The per-dimension gate mechanism} provides
inherently more capacity than scalar gating, but raises the risk of overfitting on
small clinical datasets---the ablation comparing per-dimension vs. scalar gates
(ablation~\#2) is designed to quantify this trade-off. \textbf{The session-level
Transformer} captures full-session context at $O(T^2)$ cost; for $T = 512$ utterances,
this produces approximately 262k attention scores per layer, which is manageable on a
single A100 but may require linear attention variants for edge deployment.

Compared to existing clinical MER architectures, our design uniquely integrates four
capabilities: adaptive cross-modal gating informed by learned modality quality,
session-level temporal modelling with therapeutic phase awareness, clinically grounded
interpretability via a concept bottleneck, and parameter-efficient clinical adaptation
via LoRA. No prior single architecture combines all four.

\paragraph{Limitations.}
The primary limitation is that no model has been trained or evaluated on real clinical
data. All results reported are from the validation harness operating on synthetic data
with randomly initialised weights. Consequently:

\begin{itemize}
  \item All performance targets from the research contract (MAE $< 3.0$ on E-DAIC,
        RMSE $< 3.5$, $<15\%$ degradation under 40\% dropout, $\geq 10\%$ sustained
        affect F1 gain) remain \texttt{TODO: unverified}.
  \item The quality estimators require supervised training on synthetic corruptions
        before their sensitivity can be assessed; random-initialisation benchmarks
        show low sensitivity as expected.
  \item Phase labelling relies on heuristic time-window assignment; inter-annotator
        agreement with clinical experts has not been measured.
  \item The clinically-correlated dropout simulation protocol (three patterns) is
        heuristic, not validated against real clinical dropout distributions.
  \item Fairness evaluation and clinician interpretability studies are deferred pending
        access to sufficiently diverse clinical data and domain experts.
  \item The MDD-MARF baseline~\cite{zhou2025mddmarf} (current SOTA on DAIC-WOZ) has
        not been re-implemented for comparison.
\end{itemize}


\section{Conclusion}
\label{sec:conclusion}

We have presented a reference architecture and comprehensive validation harness for
multimodal emotion recognition in mental health diagnostics. The architecture introduces
Cross-Modal Gated Exchange Fusion (CM-GEF) with per-dimension feature substitution
guided by learned modality quality estimates, a Session-Level Hierarchical Transformer
(SLHT) with therapeutic phase embeddings that model the temporal structure of clinical
interviews, and a Clinical Concept Bottleneck (CCB) that maps session representations
through DSM-5-aligned clinical constructs for interpretable PHQ-8 prediction. The
validation harness provides 35+ unit tests, seven synthetic benchmarks, and ten
single-field ablations, all operating on synthetic data with no dependency on clinical
datasets.

We release the typed \texttt{ModelConfig} contract, the complete PyTorch reference
implementation, and the validation suite as an open-source foundation for reproducible
clinical MER research. Our contribution is a design artifact---the architecture and the
infrastructure to characterise it---not a trained model with empirical superiority claims.

Concrete follow-up directions include: (1) integrating DAIC-WOZ and E-DAIC data loaders
and conducting the two-stage training pipeline to measure real clinical metrics;
(2) implementing the MDD-MARF baseline and comparing against CM-GEF on identical
clinical train/test splits; (3) running systematic label-efficiency sweeps ($5\%$,
$10\%$, $25\%$, $50\%$, $100\%$ of clinical data) to validate the LoRA adaptation
hypothesis; (4) conducting a clinician Likert-scale study to assess the interpretability
of concept-bottleneck explanations; and (5) evaluating demographic fairness on
population subgroups where data permits. We hope this architecture and infrastructure
serve the community as a reproducible starting point for clinically-validated MER.


% Acknowledgments omitted for anonymous submission.
% \begin{ack}
% This work was supported by ...
% \end{ack}


\section*{References}

{\small
\begin{thebibliography}{99}

\bibitem{who2023depression}
  World Health Organization.\ (2023).
  Depressive disorder (depression).
  \textit{WHO Fact Sheet}.

\bibitem{survey2026multimodal}
  Survey: State-of-the-art Multimodal Emotion Recognition (2026).
  \textit{Intelligent Systems with Applications}.

\bibitem{chen2022wavlm}
  Chen, S., Wang, C., Chen, Z., Wu, Y., Liu, S., Chen, Z., Li, J., Kanda, N.,
  Yoshioka, T., Xiao, X., \& others.\ (2022).
  WavLM: Large-scale self-supervised pre-training for full stack speech processing.
  \textit{IEEE Journal of Selected Topics in Signal Processing}, 16(6), 1505--1518.

\bibitem{hsu2021hubert}
  Hsu, W.-N., Bolte, B., Tsai, Y.-H.H., Lakhotia, K., Salakhutdinov, R., \&
  Mohamed, A.\ (2021).
  HuBERT: Self-supervised speech representation learning by masked prediction of hidden
  units. \textit{IEEE/ACM Transactions on Audio, Speech, and Language Processing}, 29,
  3451--3460.

\bibitem{oquab2023dinov2}
  Oquab, M., Darcet, T., Moutakanni, T., Vo, H., Szafraniec, M., Khalidov, V.,
  Fernandez, P., Haziza, D., Massa, F., El-Nouby, A., \& others.\ (2023).
  DINOv2: Learning robust visual features without supervision.
  \textit{arXiv:2304.07193}.

\bibitem{ji2022mentalbert}
  Ji, S., Zhang, T., Ansari, L., Fu, J., Tiwari, P., \& Cambria, E.\ (2022).
  MentalBERT: Publicly available readable language model for mental health.
  \textit{Proceedings of the 13th International Workshop on Health Text Mining and
  Information Analysis}, 64--72.

\bibitem{tsai2019mult}
  Tsai, Y.-H.H., Bai, S., Liang, P.P., Kolter, J.Z., Morency, L.-P., \&
  Salakhutdinov, R.\ (2019).
  Multimodal Transformer for unaligned multimodal language sequences.
  \textit{Proceedings of the 57th Annual Meeting of the ACL}, 6558--6569.

\bibitem{aefnet2025}
  AEFNet (2025).
  Adaptive External Attention-Enhanced Fusion Network for multimodal emotion recognition.
  \textit{Alexandria Engineering Journal}.

\bibitem{amte2025}
  AMTE (2025).
  Adaptive Multimodal Transformer based on Exchanging for multimodal emotion recognition.
  \textit{Scientific Reports}.

\bibitem{micga2025}
  Mi-CGA (2025).
  Cross-modal Graph Attention for incomplete multimodal emotion recognition.
  \textit{Neurocomputing}.

\bibitem{unifiedmer2025}
  Unified MER Framework (2025).
  Homogeneous and heterogeneous multimodal fusion for emotion recognition.
  \textit{Information Fusion}.

\bibitem{sadeghi2024}
  Sadeghi, M., \etal\ (2024).
  Pr3+Whisper+AudioQual for depression detection on E-DAIC.
  \textit{Interspeech / IEEE}.

\bibitem{zhou2025mddmarf}
  Zhou, Ge, \etal\ (2025).
  MDD-MARF: Multi-level attention residual fusion for depression detection.
  \textit{Journal of Biomedical Informatics}.

\bibitem{dstc2025}
  DSTC (2025).
  Temporal ProbSparse Attention Transformer for depression detection from multimodal data.
  \textit{IEEE Transactions on Affective Computing}.

\bibitem{dsin2025}
  DSIN (2025).
  Deep Spatiotemporal Interaction Network for multimodal emotion recognition.
  \textit{Information Sciences}.

\bibitem{ecmc2026}
  ECMC (2026).
  Emotion-Cognition Captioning for Mental Health.
  \textit{Proceedings of AAAI 2026}.

\bibitem{koh2020concept}
  Koh, P.W., Nguyen, T., Tang, Y.S., Mussmann, S., Pierson, E., Kim, B., \&
  Liang, P.\ (2020).
  Concept bottleneck models.
  \textit{Proceedings of ICML 2020}, 5338--5348.

\bibitem{kroenke2009phq}
  Kroenke, K., Strine, T.W., Spitzer, R.L., Williams, J.B.W., Berry, J.T., \&
  Mokdad, A.H.\ (2009).
  The PHQ-8 as a measure of current depression in the general population.
  \textit{Journal of Affective Disorders}, 114(1--3), 163--173.

\bibitem{survey2025mer}
  Survey: Recent Advances in Multimodal Affective Computing (2026).
  \textit{arXiv:2409.07388}.

\end{thebibliography}
}


\newpage
\section*{NeurIPS Paper Checklist}

\paragraph{1. Claims.}
Do the main claims made in the abstract and introduction accurately reflect the paper's
contributions and scope?
\textbf{Yes.}
The claims are scoped to a reference architecture and validation harness, not to
trained-model performance on clinical data.

\paragraph{2. Experiments.}
Does the paper describe all experiments and are they reproducible?
\textbf{Yes.}
All experiments are synthetic benchmarks running via a single \texttt{pytest} or
\texttt{python} command. Commands are documented in each source file header.

\paragraph{3. Theory.}
Does the paper include theoretical results?
\textbf{No.}
The paper presents an architecture with inductive-bias justifications and complexity
analysis but no formal theoretical bounds.

\paragraph{4. Datasets.}
Does the paper discuss datasets used?
\textbf{No.}
All experiments use synthetic data. Real clinical datasets (DAIC-WOZ, E-DAIC) are
discussed as intended targets for future work but are not used in this paper.

\paragraph{5. Code.}
Does the paper provide code and is it well-documented?
\textbf{Yes.}
The full PyTorch reference implementation, typed configuration, and validation suite
are released as open source. All files include docstrings and shape annotations.

\paragraph{6. Broader Impacts.}
Does the paper discuss potential broader impacts?
\textbf{Yes.}
The paper discusses the intended clinical application, the importance of demographic
fairness evaluation, and the need for clinician-in-the-loop validation before deployment.
Potential risks of deployment without adequate validation are acknowledged in the
limitations section.

\paragraph{7. Compute.}
Does the paper report compute resources used?
\textbf{Partial.}
The validation harness runs on CPU (single node, 8 cores). bf16 numerical stability
tests require an NVIDIA A100 80~GB GPU. Training resource estimates (e.g., ``4 hours
on 4$\times$A100'' for Stage~1) are documented as intended configurations in the
training recipe but have not been empirically verified.

\end{document}
